% !TITLE 南陵别儿童入京
% !AUTHOR 李白
% !DYNASTY 唐
% !ORDER 8

\begin{poem}
白酒新熟山中归，黄鸡啄黍\textsuperscript{1}秋正肥。
呼童烹鸡酌白酒，儿女嬉笑牵人衣。
高歌取醉欲自慰，起舞落日争光辉。
游说万乘\textsuperscript{2}苦不早，著鞭跨马涉远道。
会稽愚妇轻买臣\textsuperscript{3}，余亦辞家西入秦\textsuperscript{4}。
仰天大笑出门去，我辈岂是蓬蒿人\textsuperscript{5}。
\end{poem}

\begin{annotations}
\notetext{1}{黍：黄米，一种谷物。秋天黄鸡啄食黍米，长得正肥。写农家丰收的景象。}
\notetext{2}{游说万乘：游说（shuì），向君主进言献策；万乘（shèng），一万辆兵车，代指皇帝。李白遗憾没能早点去游说皇帝。}
\notetext{3}{会稽愚妇轻买臣：西汉朱买臣家贫好读书，妻子嫌他穷而离去，后来朱买臣做了会稽太守。李白借用此典暗讽那些看不起自己的人。}
\notetext{4}{西入秦：向西进入秦地（长安在陕西，春秋战国属秦）。李白接到皇帝诏书，辞家西去长安。}
\notetext{5}{蓬蒿人：住在蓬草和蒿草之间的草野之人，指平民百姓。我辈岂是普通人！}
\end{annotations}

\begin{appreciation}
《南陵别儿童入京》写的是李白人生里最得意的一天。天宝元年（742年），四十二岁的李白终于接到唐玄宗的诏书，召他进京。诗题里的“儿童”是他的孩子，“入京”是去长安。离开家之前，他写了这首诗，全家都听见了他的笑声。

诗从农家秋景写起：“白酒新熟山中归，黄鸡啄黍秋正肥。”酒刚酿好，黄鸡啄着黍米，秋天正肥。诏书到来之前，他只是南陵乡下的一个中年人。“呼童烹鸡酌白酒，儿女嬉笑牵人衣。”叫童子杀鸡温酒，孩子们笑着拉住他的衣角。

“高歌取醉欲自慰，起舞落日争光辉。”他高歌、痛饮、起舞。这里有个词值得停下来：自慰。自我安慰。得意到了需要自我安慰的地步——那说明前面二十年确实苦。二十年的漫游、求见、碰壁，他从没在人前说过一句软话，可这一句“自慰”漏了底：他也需要为自己高兴一下。

“会稽愚妇轻买臣，余亦辞家西入秦。”用了西汉朱买臣的典故。朱买臣家贫，靠砍柴读书，妻子嫌他穷，离开了他；后来朱买臣做了会稽太守。李白说：当年瞧不起我的人，就像那个愚妇。这不是大度，是积了二十年的委屈，终于找到了出口。

“仰天大笑出门去，我辈岂是蓬蒿人。”蓬蒿人，就是草野间的人。我李白，怎么可能是平凡人。这是全诗最响的一句。还记得《上李邕》吗？二十年前那个说自己是大鹏的少年，终于等到了风。同一只大鹏，同一种“我辈”。

后来的事你也知道：他在长安只待了不到三年就被赐金放还——这本书后面那首《行路难》，写的就是那之后的事。但“仰天大笑出门去”这一声笑是真的。你将来考出了好成绩，尽可以笑得大声些——那是你以前攒下的委屈换来的，你配得上。
\end{appreciation}
